\section{Actual affine covers recorded by finite Rees trees}

\begin{definition}[Leaf occurrences]\label{def:sc0-rees-leaf-occurrences}
\lean{EnvelopeProperty.ReesChartTree.LeafIndex,EnvelopeProperty.ReesChartTree.leafAlgebra}
\leanok
A leaf occurrence of a finite Rees tree records all branch indices on its path.
The one-leaf tree has one occurrence. The occurrences at a node are the disjoint
union, over its nonzero generator indices, of the occurrences of its children.
Each occurrence has its corresponding leaf subalgebra; different occurrences may
have the same subalgebra.
\end{definition}

\begin{lemma}[Finitely many actual leaf occurrences]\label{lem:sc0-finite-rees-leaf-occurrences}
\lean{EnvelopeProperty.ReesChartTree.finite_leafIndex,EnvelopeProperty.ReesChartTree.nonempty_leafIndex,EnvelopeProperty.ReesChartTree.leafAlgebra_mem_leaves,EnvelopeProperty.ReesChartTree.exists_leafIndex_of_mem_leaves}
\leanok
\uses{def:sc0-rees-leaf-occurrences}
The occurrences of a finite Rees tree form a finite nonempty set, and their
corresponding subalgebras are exactly the leaf subalgebras of the tree.
\end{lemma}
\begin{proof}
Induct on the tree. The assertion for a one-leaf tree follows from its single
occurrence. At a node the generator index set is finite and nonempty by the
nonzero-generator condition. Each child has finitely many nonempty occurrences
by induction, so their disjoint union is finite and nonempty. The set of
corresponding algebras is the union of the children's leaf sets, which is exactly
the leaf set of the node.
\end{proof}

\begin{theorem}[An actual blowup sequence gives a finite Rees cover]\label{thm:sc0-actual-sequence-rees-cover}
\lean{EnvelopeProperty.Char0.exists_actual_rees_cover,EnvelopeProperty.Char0.reesLeafAffineOpenCover,EnvelopeProperty.Char0.finite_reesLeafAffineOpenCover_index}
\leanok
\uses{lem:sc0-finite-rees-leaf-occurrences,lem:sc0-affine-centre-finite-generators,lem:sc0-global-affine-rees-cover,lem:sc0-rees-tail-open-iso,lem:sc0-sequence-last-pullback,lem:sc0-open-iso-source-point}
Let $B$ be Noetherian, let $K$ be a field with a $B$-algebra structure, and let
$C\subseteq K$ be a finitely generated $B$-subalgebra. An actual finite blowup
sequence over $\operatorname{Spec}C$ which is an isomorphism over a nonempty open
admits a finite Rees tree rooted at $C$ and an actual affine open cover of its
final scheme indexed by the tree's leaf occurrences. The coordinate ring at each
occurrence is its literal leaf subalgebra of $K$, and its composite to
$\operatorname{Spec}C$ is induced by the original algebra inclusion.
\end{theorem}
\begin{proof}
Induct on the sequence length, simultaneously for all such finitely generated
subalgebras. The empty sequence gives the one-leaf tree and identity affine cover.
For a nonempty sequence, its nonempty isomorphism open supplies a point in its
final scheme. Project this point through the remaining sequence to obtain a point
of the first blowup. Its centre therefore has a finite generating family with
a nonzero member. The nonzero ratio charts of that first blowup form its actual
finite affine open cover. Their rings are finitely generated $B$-subalgebras of
$K$, hence Noetherian. Pull the tail sequence back to each chart. Its length is
strictly smaller, and Lemma~\ref{lem:sc0-rees-tail-open-iso} gives it a nonempty
isomorphism open. Apply induction and attach the resulting child trees at the
first node. Compose the leaf open immersions with the actual final-stage lifts.
These are open immersions by the cartesian final-stage squares. They cover:
given a final point, its projection to the first blowup lies in one ratio chart;
the cartesian square lifts that point to the pulled-back final scheme, where the
inductive cover supplies a leaf occurrence. The cartesian commutativity and the
ratio-chart projection identities give the asserted composite inclusion maps.
\end{proof}
